% Options for packages loaded elsewhere
\PassOptionsToPackage{unicode}{hyperref}
\PassOptionsToPackage{hyphens}{url}
\documentclass[
]{article}
\usepackage{xcolor}
\usepackage{amsmath,amssymb}
\setcounter{secnumdepth}{-\maxdimen} % remove section numbering
\usepackage{iftex}
\ifPDFTeX
  \usepackage[T1]{fontenc}
  \usepackage[utf8]{inputenc}
  \usepackage{textcomp} % provide euro and other symbols
\else % if luatex or xetex
  \usepackage{unicode-math} % this also loads fontspec
  \defaultfontfeatures{Scale=MatchLowercase}
  \defaultfontfeatures[\rmfamily]{Ligatures=TeX,Scale=1}
\fi
\usepackage{lmodern}
\ifPDFTeX\else
  % xetex/luatex font selection
\fi
% Use upquote if available, for straight quotes in verbatim environments
\IfFileExists{upquote.sty}{\usepackage{upquote}}{}
\IfFileExists{microtype.sty}{% use microtype if available
  \usepackage[]{microtype}
  \UseMicrotypeSet[protrusion]{basicmath} % disable protrusion for tt fonts
}{}
\makeatletter
\@ifundefined{KOMAClassName}{% if non-KOMA class
  \IfFileExists{parskip.sty}{%
    \usepackage{parskip}
  }{% else
    \setlength{\parindent}{0pt}
    \setlength{\parskip}{6pt plus 2pt minus 1pt}}
}{% if KOMA class
  \KOMAoptions{parskip=half}}
\makeatother
\usepackage{longtable,booktabs,array}
\usepackage{calc} % for calculating minipage widths
% Correct order of tables after \paragraph or \subparagraph
\usepackage{etoolbox}
\makeatletter
\patchcmd\longtable{\par}{\if@noskipsec\mbox{}\fi\par}{}{}
\makeatother
% Allow footnotes in longtable head/foot
\IfFileExists{footnotehyper.sty}{\usepackage{footnotehyper}}{\usepackage{footnote}}
\makesavenoteenv{longtable}
\setlength{\emergencystretch}{3em} % prevent overfull lines
\providecommand{\tightlist}{%
  \setlength{\itemsep}{0pt}\setlength{\parskip}{0pt}}
\usepackage{bookmark}
\IfFileExists{xurl.sty}{\usepackage{xurl}}{} % add URL line breaks if available
\urlstyle{same}
\hypersetup{
  hidelinks,
  pdfcreator={LaTeX via pandoc}}

\author{}
\date{}

\begin{document}

\textbf{Оглавление}

\begin{enumerate}
\def\labelenumi{\arabic{enumi}.}
\item
  \textbf{Титульный лист}
\item
  \textbf{Введение}

  \begin{itemize}
  \item
    \textbf{2.1. Актуальность проекта}
  \item
    \textbf{2.2. Цель проекта}
  \item
    \textbf{2.3. Задачи проекта}
  \end{itemize}
\item
  \textbf{Общее описание игры}
\item
  \textbf{Игровой процесс}

  \begin{itemize}
  \item
    \textbf{4.1. Ключевые особенности геймплея}
  \item
    \textbf{4.2. Основной игровой цикл}
  \end{itemize}
\item
  \textbf{Сюжет и игровой прогресс}

  \begin{itemize}
  \item
    \textbf{5.1. Завязка}
  \item
    \textbf{5.2. Прогрессия}
  \item
    \textbf{5.3. Условие победы}
  \end{itemize}
\item
  \textbf{Техническая реализация}

  \begin{itemize}
  \item
    \textbf{6.1. Движок}
  \item
    \textbf{6.2. Основные системы}
  \item
    \textbf{6.3. Планируемые инструменты}
  \end{itemize}
\item
  \textbf{Этапы разработки}
\item
  \textbf{Ожидаемые результаты}
\item
  \textbf{Заключение}
\end{enumerate}

\textbf{Документация проекта: 2D Top-Down Action}

\begin{center}\rule{0.5\linewidth}{0.5pt}\end{center}

\textbf{1. Титульный лист}

\textbf{Название проекта:}~2D Top-Down Action\\
\textbf{Автор:}~ученик 10 класса\\
\textbf{Платформа:}~Unity\\
\textbf{Жанр:}~2D экшен с элементами исследования (вид сверху)\\
\textbf{Год:}~2026

\begin{center}\rule{0.5\linewidth}{0.5pt}\end{center}

\textbf{2. Введение}

\textbf{2.1. Актуальность проекта}

Современные инди-игры всё чаще обращаются к атмосферным проектам, где
ключевую роль играют свет и темнота. Механика освещения как центральный
элемент геймплея позволяет создать уникальный опыт и выделить игру среди
аналогов. Проект сочетает динамичный экшен с исследованием, что делает
его интересным для широкой аудитории.

\textbf{2.2. Цель проекта}

Создать 2D экшен с видом сверху, в котором игрок проходит ряд испытаний,
преодолевает сопротивление противников и достигает финальной точки
уровня для успешного завершения игры.

\textbf{2.3. Задачи проекта}

\begin{enumerate}
\def\labelenumi{\arabic{enumi}.}
\item
  Разработать концепцию игрового мира и сюжета.
\item
  Реализовать механику освещения как центральный элемент геймплея.
\item
  Создать систему взаимодействия с NPC.
\item
  Реализовать систему сбора деталей и их установки на вышку.
\item
  Разработать боевую систему и противников.
\item
  Реализовать финальную сцену и условие победы.
\end{enumerate}

\begin{center}\rule{0.5\linewidth}{0.5pt}\end{center}

\textbf{3. Общее описание игры}

\begin{longtable}[]{@{}
  >{\raggedright\arraybackslash}p{(\linewidth - 2\tabcolsep) * \real{0.3536}}
  >{\raggedright\arraybackslash}p{(\linewidth - 2\tabcolsep) * \real{0.6464}}@{}}
\toprule\noalign{}
\begin{minipage}[b]{\linewidth}\raggedright
Параметр
\end{minipage} & \begin{minipage}[b]{\linewidth}\raggedright
Значение
\end{minipage} \\
\midrule\noalign{}
\endhead
\bottomrule\noalign{}
\endlastfoot
Жанр & 2D экшен с элементами исследования \\
Платформа & Unity(v6) \\
Вид камеры & Вид сверху \\
Визуальный стиль & Тёмная карта, контрастное освещение \\
Ключевая механика & Освещение \\
\end{longtable}

\begin{center}\rule{0.5\linewidth}{0.5pt}\end{center}

\textbf{4. Игровой процесс}

\textbf{4.1. Ключевые особенности геймплея}

\textbf{Визуальный стиль и освещение:}

\begin{itemize}
\item
  Карта полностью погружена в темноту.
\item
  Источниками света выступают исключительно сам персонаж и враги.
\item
  Механика освещения является центральным элементом геймплея.
\item
  Перспектива --- вид сверху.
\end{itemize}

\textbf{Взаимодействие с нейтральными персонажами (NPC):}

\begin{itemize}
\item
  Игрок может взаимодействовать с NPC для получения информации,
  подсказок или квестов.
\item
  NPC могут помогать в исследовании или предоставлять доступ к новым
  зонам.
\end{itemize}

\textbf{4.2. Основной игровой цикл}

\begin{enumerate}
\def\labelenumi{\arabic{enumi}.}
\item
  Перемещение по тёмной карте.
\item
  Взаимодействие с NPC.
\item
  Поиск деталей.
\item
  Бой с врагами (преодоление испытаний).
\item
  Установка деталей на вышку.
\item
  Финал игры.
\end{enumerate}

\begin{center}\rule{0.5\linewidth}{0.5pt}\end{center}

\textbf{5. Сюжет и игровой прогресс}

\textbf{5.1. Завязка}

Персонаж оказывается в незнакомой локации, погружённой во тьму.
Единственный источник света --- он сам и враги. Чтобы выбраться, ему
необходимо восстановить работу электрической вышки.

\textbf{5.2. Прогрессия}

Для завершения игры и выхода из локации игроку необходимо:

\begin{enumerate}
\def\labelenumi{\arabic{enumi}.}
\item
  \textbf{Исследовать территорию}~и взаимодействовать с другими
  персонажами.
\item
  \textbf{Собрать все необходимые детали (компоненты).}
\item
  \textbf{Активировать электрическую вышку.}
\end{enumerate}

\textbf{5.3. Условие победы}

Активация вышки запускает финальную сцену и завершает игровой сеанс.

\begin{center}\rule{0.5\linewidth}{0.5pt}\end{center}

\textbf{6. Техническая реализация}

\textbf{6.1. Движок}

\begin{itemize}
\item
  \textbf{Unity}~--- (Unity v6) выбран благодаря мощным инструментам для
  работы с 2D-графикой и освещением (2D Lighting System).
\end{itemize}

\textbf{6.2. Основные системы}

\begin{longtable}[]{@{}
  >{\raggedright\arraybackslash}p{(\linewidth - 2\tabcolsep) * \real{0.3551}}
  >{\raggedright\arraybackslash}p{(\linewidth - 2\tabcolsep) * \real{0.6449}}@{}}
\toprule\noalign{}
\begin{minipage}[b]{\linewidth}\raggedright
Система
\end{minipage} & \begin{minipage}[b]{\linewidth}\raggedright
Описание
\end{minipage} \\
\midrule\noalign{}
\endhead
\bottomrule\noalign{}
\endlastfoot
Система освещения & Динамическое освещение от персонажа и врагов \\
Система перемещения & 2D-движение с видом сверху \\
Система взаимодействия & Взаимодействие с NPC и объектами \\
Система сбора предметов & Поиск и подбор деталей \\
Боевая система & Бой с врагами \\
Система прогресса & Установка деталей на вышку, активация финала \\
\end{longtable}

\textbf{6.3. Планируемые инструменты}

\begin{itemize}
\item
  Unity v6 (2D)
\item
  Lighting
\item
  Tilemap
\item
  Animator
\item
  Physics 2D
\item
  UI System
\item
  Visual Studio
\item
  Искусственный интеллект(для создания карты, анимаций и персонажей)
\end{itemize}

\begin{center}\rule{0.5\linewidth}{0.5pt}\end{center}

\textbf{7. Этапы разработки}

\begin{longtable}[]{@{}
  >{\raggedright\arraybackslash}p{(\linewidth - 4\tabcolsep) * \real{0.1145}}
  >{\raggedright\arraybackslash}p{(\linewidth - 4\tabcolsep) * \real{0.5265}}
  >{\raggedright\arraybackslash}p{(\linewidth - 4\tabcolsep) * \real{0.3590}}@{}}
\toprule\noalign{}
\begin{minipage}[b]{\linewidth}\raggedright
Этап
\end{minipage} & \begin{minipage}[b]{\linewidth}\raggedright
Содержание
\end{minipage} & \begin{minipage}[b]{\linewidth}\raggedright
Срок
\end{minipage} \\
\midrule\noalign{}
\endhead
\bottomrule\noalign{}
\endlastfoot
1 & Проектирование концепции и сюжета & 1 неделя \\
2 & Создание прототипа уровня & 2 недели \\
3 & Реализация системы освещения & 2 недели \\
4 & Добавление NPC и взаимодействий & 1 неделя \\
5 & Реализация боевой системы & 2 недели \\
6 & Система сбора деталей и вышка & 1 неделя \\
7 & Финальная сцена и полировка & 1 неделя \\
8 & Тестирование и отладка & 1 неделя \\
\end{longtable}

\begin{center}\rule{0.5\linewidth}{0.5pt}\end{center}

\textbf{8. Ожидаемые результаты}

По завершении проекта будет получена игра, в которой:

\begin{itemize}
\item
  Реализована уникальная механика освещения.
\item
  Присутствует полноценный игровой цикл.
\item
  Работает система взаимодействия с NPC.
\item
  Реализован сбор деталей и активация вышки.
\item
  Игра завершается финальной сценой.
\end{itemize}

\begin{center}\rule{0.5\linewidth}{0.5pt}\end{center}

\textbf{9. Заключение}

Проект «2D Top-Down Action» представляет собой атмосферный экшен с видом
сверху, где темнота и свет являются не просто визуальным решением, а
основой геймплея. Реализация проекта позволит освоить ключевые навыки
работы с Unity, 2D-графикой, системами освещения и игровым дизайном.
Игра может быть интересна как для портфолио, так и для дальнейшего
развития в полноценный коммерческий продукт.

\end{document}
